\documentclass[12pt, titlepage]{article}

\usepackage{amsmath, mathtools}

\usepackage[round]{natbib}
\usepackage{amsfonts}
\usepackage{amssymb}
\usepackage{graphicx}
\usepackage{colortbl}
\usepackage{xr}
\usepackage{hyperref}
\usepackage{longtable}
\usepackage{xfrac}
\usepackage{tabularx}
\usepackage{float}
\usepackage{siunitx}
\usepackage{booktabs}
\usepackage{multirow}
\usepackage[section]{placeins}
\usepackage{caption}
\usepackage{fullpage}

\hypersetup{
bookmarks=true,     % show bookmarks bar?
colorlinks=true,       % false: boxed links; true: colored links
linkcolor=red,          % color of internal links (change box color with linkbordercolor)
citecolor=blue,      % color of links to bibliography
filecolor=magenta,  % color of file links
urlcolor=cyan          % color of external links
}

\usepackage{array}

\externaldocument{../../SRS/SRS}

\input{../../Comments.text}
\input{../../Common.text}
\setlength{\parindent}{0pt}
\setlength{\parskip}{0.75em}
\setlength{\emergencystretch}{3em}
\newcommand{\True}{\textnormal{true}}
\newcommand{\False}{\textnormal{false}}
\newcommand{\nul}{\textnormal{null}}
\newcommand{\Mref}[1]{\hyperref[sec:M#1]{M#1}}
\newcommand{\noinit}{$\neg$initialized $\Rightarrow$ NotInitialized}
\begin{document}

\title{Module Interface Specification for \progname{}}

\author{\authname}

\date{\today}

\maketitle

\pagenumbering{roman}

\section{Revision History}

\begin{tabularx}{\textwidth}{p{3cm}p{2cm}X}
\toprule {\bf Date} & {\bf Version} & {\bf Notes}\\
\midrule
Date 1 & 1.0 & Notes\\
Date 2 & 1.1 & Notes\\
\bottomrule
\end{tabularx}

~\newpage

\section{Symbols, Abbreviations and Acronyms}

See SRS Documentation at \url{https://softwaremandems.github.io/FPVR/SRS/SRS.pdf}
in the Glossary section for the full table of symbols, abbreviations and acronyms.

% \wss{Also add any additional symbols, abbreviations or acronyms}

\newpage

\tableofcontents

\newpage

\pagenumbering{arabic}

\section{Introduction}

The following document details the Module Interface Specifications for
FPVR: An FPV drone simulation game. The application is designed to provide a realistic
virtual environment for pilots to hone their skills by completing increasingly difficult
courses. FPVR is powered by the Unity Game Engine, connected to a backend server and database to
manage user accounts and scores, hardware controller compatible, and features an AI-assisted interface
to generate feedback and suggestions for improvement.

Complementary documents include the System Requirement Specifications
and Module Guide.  The full documentation and implementation can be
found at \url{https://github.com/SoftwareMandems/FPVR}.

\section{Notation}

% \wss{You should describe your notation.  You can use what is below as
%   a starting point.}

The structure of the MIS for modules comes from \citet{HoffmanAndStrooper1995},
with the addition that template modules have been adapted from
\cite{GhezziEtAl2003}.  The mathematical notation comes from Chapter 3 of
\citet{HoffmanAndStrooper1995}.  For instance, the symbol := is used for a
multiple assignment statement and conditional rules follow the form $(c_1
\Rightarrow r_1 | c_2 \Rightarrow r_2 | ... | c_n \Rightarrow r_n )$.

The following table summarizes the primitive data types used by \progname. 

\begin{center}
\renewcommand{\arraystretch}{1.2}
\noindent 
\begin{tabular}{l l p{7.5cm}} 
\toprule 
\textbf{Data Type} & \textbf{Notation} & \textbf{Description}\\ 
\midrule
character & char & a single symbol or digit\\
integer & $\mathbb{Z}$ & a number without a fractional component in (-$\infty$, $\infty$) \\
natural number & $\mathbb{N}$ & a number without a fractional component in [1, $\infty$) \\
real & $\mathbb{R}$ & any number in (-$\infty$, $\infty$)\\
boolean & $\mathbb{B}$ & a value that is either \True{} or \False{}\\
string & $\Sigma^*$ & a sequence of characters (char)\\
enumerate & $\{\text{E}_i \:|\: i \leq n, \text{E}_i \in \Sigma^*\}$ & a finite unordered list of named constants\\
sequence & seq of T & an ordered list of elements of type T\\
tuple & (T$_1$, T$_2$, ..., T$_n$) & a finite ordered list of values of types T$_1$, T$_2$, ..., T$_n$\\
dictionary & dict of T$_1$ to T$_2$ & a key-value pairing from type T$_1$ to type T$_2$\\
function & T$_1$ $\rightarrow$ T$_2$ & a mapping from type T$_1$ to type T$_2$\\
set & set of T & an unordered collection of distinct elements of type T\\
\bottomrule
\end{tabular}
\end{center}

\noindent
The value \nul{} marks a field or variable that has no value, for example no
course selected yet. A type that allows it is written T $\cup \ \{\nul\}$. For a
tuple $x$ with a field named $f$, $x.f$ is the value of that field. For a
sequence $s$, $|s|$ is its length, $s[i]$ is its $i$-th element (starting at
1), $s \,||\, t$ is $s$ followed by $t$, and $\langle\rangle$ is the empty
sequence. A multiple assignment $v_1, v_2 := e_1, e_2$ evaluates every
expression first and then assigns them together. Following the MIS format, a
module heading names its kind: Module (a library or abstract object) or
Template Module (an abstract data type).

\subsection*{Shared Types}

Two types are fixed by the SRS rather than by any module's design, and modules
that share no uses relation both need them. They are defined once here so that
each module's Uses list can match the MG uses hierarchy exactly.

\begin{itemize}
\item Pose = tuple of (position: $\mathbb{R}^3$, orientation: $\mathbb{R}^3$).
  Position is in metres in the course frame, with $\hat{z} = (0, 0, 1)$
  pointing up. Orientation is the roll $\phi$, pitch $\theta$ and yaw $\psi$
  angles in radians. Used for the drone, the headset, course starts and
  checkpoints (FR-FLT-01, FR-CHK-05).
\item Rating = $\{$GOLD, SILVER, BRONZE, NEEDS\_IMPROVEMENT$\}$. The four
  rating names required by FR-PER-09. Which scores map to which rating is the
  secret of \Mref{14}.
\end{itemize}

\subsection*{Exported Types}

Every other type is exported by the module that owns it (Table~\ref{TblTypes}).
A module that needs one of these types lists its owner under Uses.

\begin{table}[h!]
\centering
\renewcommand{\arraystretch}{1.15}
\begin{tabularx}{\textwidth}{>{\raggedright\arraybackslash}p{5.6cm} l >{\raggedright\arraybackslash}X}
\toprule
\textbf{Types} & \textbf{Owner} & \textbf{Meaning}\\
\midrule
ExecutionMode, View & \Mref{1} & Standalone or PCVR; an image shown in the headset\\
DeviceType, ControlInput, MenuInput, ControlMapping & \Mref{2} & The active controller, the four flight axes, a menu button and the axis bindings\\
ActionKind, UserAction, Screen, CourseItem, HudView, ReviewView, BoardRow, HistoryRow & \Mref{3} & A user's choice and the data each screen displays\\
Phase & \Mref{4} & The stage of the current training session\\
Mode, Difficulty, CourseInfo & \Mref{5} & Training mode and course description\\
Passage & \Mref{6} & A checkpoint gate or the finish the drone passed through\\
Evaluation & \Mref{8} & The frozen measurements, score and rating of an attempt\\
User & \Mref{9} & A signed-in player and their session\\
Scope, ResultInput, LeaderboardEntry, HistoryEntry & \Mref{10} & Result data passed in and out of synchronization\\
FlightState & \Mref{11} & The drone's pose, velocity and angular velocity\\
SyncStatus, ResultRecord & \Mref{12} & A stored result and its synchronization state\\
AttemptSummary, Feedback & \Mref{13} & The input and output of feedback analysis\\
RankKey & \Mref{15} & The values used to order one result\\
\bottomrule
\end{tabularx}
\caption{Exported types and the modules that own them}
\label{TblTypes}
\end{table}

\section{Module Decomposition}

Table~\ref{TblMH} reproduces the module hierarchy from the Module Guide (MG).
The modules are grouped by the kind of secret they hide. Hardware-hiding modules
hide the headset and controller devices. Behaviour-hiding modules hide the
behaviour required by the SRS. Software decision modules hide algorithms and
data structures that the SRS does not specify.

Each module keeps the number it has in the MG, so M1 to M15 refer to the same
module in both documents. Only the leaf modules in the Level 2 column are
implemented, and each one is specified in its own section of this document.
The scoring formula and rating bands (M14) and the leaderboard ordering (M15)
are separate modules, so either policy can change without changing the modules
that measure attempts (M8) or manage results (M10).

The Intelligent Feedback Analyzer (M13) is optional. FR-AI-01 to FR-AI-03
depend on a separate feasibility decision, and FR-REV-04 and FR-AI-03 require
the standard post-flight review to work whether or not M13 is available. No core
module depends on M13 being present.

\begin{table}[h!]
\centering
\begin{tabular}{p{0.3\textwidth} p{0.6\textwidth}}
\toprule
\textbf{Level 1} & \textbf{Level 2}\\
\midrule

\multirow{2}{0.3\textwidth}{Hardware-Hiding}
& \Mref{1}: VR Platform Interface\\
& \Mref{2}: Controller Input Interface\\
\midrule

\multirow{8}{0.3\textwidth}{Behaviour-Hiding}
& \Mref{3}: User Interface\\
& \Mref{4}: Training Session Manager\\
& \Mref{5}: Course and Environment Manager\\
& \Mref{6}: Flight State and Crash Manager\\
& \Mref{7}: Checkpoint and Recovery Manager\\
& \Mref{8}: Attempt Metrics Manager\\
& \Mref{9}: Account and Authentication Manager\\
& \Mref{10}: Result Synchronization Manager\\
\midrule

\multirow{5}{0.3\textwidth}{Software Decision}
& \Mref{11}: Flight Dynamics Model\\
& \Mref{12}: Data Persistence Module\\
& \Mref{13}: Intelligent Feedback Analyzer (Optional)\\
& \Mref{14}: Performance Scoring Policy\\
& \Mref{15}: Leaderboard Ordering Policy\\
\bottomrule

\end{tabular}
\caption{Module Hierarchy}
\label{TblMH}
\end{table}

\newpage

\section{MIS of \wss{Module Name}} \label{Module} \wss{Use labels for
  cross-referencing}

\wss{You can reference SRS labels, such as R\ref{R_Inputs}.}

\wss{It is also possible to use \LaTeX for hypperlinks to external documents.}

\subsection{Module}

\wss{Short name for the module}

\subsection{Uses}


\subsection{Syntax}

\subsubsection{Exported Constants}

\subsubsection{Exported Access Programs}

\begin{center}
\begin{tabular}{p{2cm} p{4cm} p{4cm} p{2cm}}
\hline
\textbf{Name} & \textbf{In} & \textbf{Out} & \textbf{Exceptions} \\
\hline
\wss{accessProg} & - & - & - \\
\hline
\end{tabular}
\end{center}

\subsection{Semantics}

\subsubsection{State Variables}

\wss{Not all modules will have state variables.  State variables give the module
  a memory.}

\subsubsection{Environment Variables}

\wss{This section is not necessary for all modules.  Its purpose is to capture
  when the module has external interaction with the environment, such as for a
  device driver, screen interface, keyboard, file, etc.}

\subsubsection{Assumptions}

\wss{Try to minimize assumptions and anticipate programmer errors via
  exceptions, but for practical purposes assumptions are sometimes appropriate.}

\subsubsection{Access Routine Semantics}

\noindent \wss{accessProg}():
\begin{itemize}
\item transition: \wss{if appropriate} 
\item output: \wss{if appropriate} 
\item exception: \wss{if appropriate} 
\end{itemize}

\wss{A module without environment variables or state variables is unlikely to
  have a state transition.  In this case a state transition can only occur if
  the module is changing the state of another module.}

\wss{Modules rarely have both a transition and an output.  In most cases you
  will have one or the other.}

\subsubsection{Local Functions}

\wss{As appropriate} \wss{These functions are for the purpose of specification.
  They are not necessarily something that is going to be implemented
  explicitly.  Even if they are implemented, they are not exported; they only
  have local scope.}

\newpage

\bibliographystyle {plainnat}
\bibliography {../../../refs/References}

\newpage

\newpage{}

\section{Appendix --- Generative AI Use Disclosure}

\wss{ChatGPT (version 5) was used to brainstorm ideas for the template for this disclose section.}

\subsection{AI Tools Used}

\wss{Give the name of the AI tool(s) used and the version.}

\subsection{How and Where Used}

\wss{Explain what the tool(s) were used for.  Examples include: brainstorming,
explaining a technical concept, reviewing the document, generating or modifying
text, generating or modifying diagrams, identifying errors or omissions or
inconsistencies.}

\wss{For each of the uses, identify which parts of the document were affected.}

\wss{You do not need to report every interaction, but you should disclose the 
uses that materially contributed to your work}

\subsection{Verification of AI-Generated Content}

\wss{\begin{itemize}
    \item Describe how the team checked AI-generated or AI-assisted material
    before including it in the document.
    \item In particular, verify factual claims, technical assertions,
    requirements, calculations, references, and other information for which an
    AI system may produce incorrect or fabricated results.
    \item \textbf{The team is responsible for the accuracy, completeness,
consistency, and appropriateness of the submitted document, regardless of
whether material was generated by a human or an AI system.}
\end{itemize}}

\newpage{}

\section{Appendix --- Reflection}

\wss{Not required for CAS 741 projects}

The information in this section will be used to evaluate the team members on the
graduate attribute of Problem Analysis and Design.

\input{../../Reflection.text}

\begin{enumerate}
  \item What went well while writing this deliverable? 
  \item What pain points did you experience during this deliverable, and how
    did you resolve them?
  \item Which of your design decisions stemmed from speaking to your client(s)
  or a proxy (e.g. your peers, stakeholders, potential users)? For those that
  were not, why, and where did they come from?
  \item While creating the design doc, what parts of your other documents (e.g.
  requirements, hazard analysis, etc), it any, needed to be changed, and why?
  \item What are the limitations of your solution?  Put another way, given
  unlimited resources, what could you do to make the project better? (LO\_ProbSolutions)
  \item Give a brief overview of other design solutions you considered.  What
  are the benefits and tradeoffs of those other designs compared with the chosen
  design?  From all the potential options, why did you select the documented design?
  (LO\_Explores)
  \item (After you have implemented another team's module, which means this
  isn't filled in until after the original deadline). What did you learn by
  implementing another team's module? Were all the details you needed in the
  documentation, or did you need to make assumptions, or ask the other team
  questions?  If your team also had another team implement one of your modules,
  what was this experience like?  Are there things in your documentation you
  could have changed to make the process go more smoothly for when an
  ``outsider'' completes some of the implementation?

  Your reflection should be about concepts learned in your SE education, such as modular decomposition, cohesion, coupling, information hiding, etc.  You should critique the quality and utility of the specification you received.

  You should create a PR for your implemented module.  A link to the PR should be part of your reflection.

\end{enumerate}

\end{document}